% Nouveau EGA V (cote 25): archive cote 25, site page 151 (PDF page 152).
% typescript. Draft compared: none.
% AI-transcribed by Claude Opus 5.5 (claude-opus-5-5, adaptive thinking, effort high), Message Batches API, batch ega25b (msgbatch_01A7n5YuXubUPNK1zsepi8RG), 2026-11-13; not yet checked by eye.
\scanpage{25/151}
% NOTE: a long diagonal pencil stroke crosses the whole page, from the lower left (near "(ii) Il existe") to the upper right (near "homomorphisme" in the first line); it runs through the text without hiding it
\mspage{12}

projectif de rang m,  $u{:}E \to I$ un homomorphisme, d'où un homomorphisme
\[ u\otimes id_{A/I} : E\otimes A/I \longrightarrow I/I^2 \ . \]
\xout{Conditions Posons} \inserted[typed]{Soit} $B = A/u(E)A$ , muni de l'augmentation $B\to B/J = A/I$
(où $J = I/u(E)A$ ). Les deux conditions suivantes sont équivalentes:

\margin{$\|$} (i) L'anneau \xout{augmenté} B \xout{\unsure{B/I est}} a une augmentation $B\to B/J$ \inserted{de type} quadra-%
tique ordinaire, de rang $n-(m-1)= n-m+1$ (le \hand{``}\uuline{rang de l'augmentation}\hand{''}
est \inserted{par définition} celui du Module projectif $J/J^2$ ).

\margin{$|$} (ii) $K = \mathrm{Ker}\ u\otimes id_{A/I}$ est un Module inversible, et l'homomor
\unsure{phisme} naturel $K \longrightarrow \mathrm{Sym}^2(J/J^2) = \mathrm{Sym}^2( (I/I^2)/\mathrm{Im}\ E\otimes A/I )$ est
``non dégénéré'' .

Il serait peut-être plus clair d'introduire une notation $V = I/I^2$ ,
et $W = J/J^2 = V / \mathrm{Im}\ E\otimes A/I$ . L'homomorphisme qualifié de naturel
dans (ii) s'obtient en notant que K s'envoie en tous cas dans
\margin{$|$} $I^2/(I^3 +$ \xout{Ixm(E)} $u(IE)) = (I^2/I^3)/\mathrm{Im}\ IE \simeq \mathrm{Sym}^2(I/I^2)/\ \mathrm{Im}\ (I/I^2)\otimes(E\otimes A/$
% NOTE: the end of this line runs off the right edge of the scan
où $I/I^2 \otimes(E \otimes A/I) \to \mathrm{Sym}^2(I/I^2)$ est défini via l'homomorphisme
canonique de multiplication $I/I^2 \otimes I/I^2 \to \mathrm{Sym}^2(I/I^2 )$ . Ces explici-%
tations posées, la démonstration de la prop.9 est à peu près triviale.

\uline{Corollaire} 10. Soit A un anneau local régulier de dimension n, $\underline{m}$ son
idéal maximal, \inserted[typed]{$k=A/\underline{m}$,} E un A-module libre de rang m, $u{:}E \to \underline{M}$ un homomorphisme,
\xout{Conditions équivalentes:} $B = A/u(E)A$, $\underline{n} = \underline{m}/u(E)A$ l'idéal maximal de B,
$V = \underline{m}/\underline{m}^2$ , $W = \underline{n}/\underline{n}^2 =$ \xout{km/} $V/\ \mathrm{Im}\ E\otimes k$ . Conditions équivalentes:

\margin{?\}} (i) B est de dimension n-m (i.e. , en termes d'une base \xout{de} $(\phi_i)$ de
\margin{$\|$} E, les\xout{t} images des $\phi_i$ dans A forment une suite A-régulière et \xout{est} le
\xout{singularité qu} point fermé de Spec(B) est une \uline{singularité quadratique}
ordinaire (par quoi nous entendons que l'augmentation $B\to B/\underline{n}$ est \inserted{de type} quadra-%
tique ordinaire).

(ii) Il existe une base $\phi_1 ,\ldots., \phi_m$ de E telle que \hand{\unsure{$\sqrt{\ }$}}$A/$\hand{$\sum_{i=1}^{m-1}$}$u(\phi i)A$

\mspage{151}
